\chapter{Hipótesis y Objetivos del Proyecto}
\label{ch:HyO}

En este capítulo se presenta la base de la investigación: la hipótesis y las 
interrogantes derivadas de la misma. El objetivo general consiste en validar 
esta premisa, para lo cual se han definido objetivos específicos que abordan 
diferentes aspectos del tema. Asimismo, se describen las tareas precisas 
requeridas para alcanzar cada objetivo y lograr resultados concluyentes. Por 
último, se exponen las limitaciones del estudio.

\section{Hipótesis}
     ``La variación sistemática de los parámetros estructurales en una arquitectura 
     reconfigurable de hardware abierto permite identificar una configuración 
     óptima que maximiza el rendimiento de temporización para un conjunto específico 
     de circuitos benchmark, demostrando que la arquitectura de la red de 
     interconexión es el factor dominante en el consumo de potencia y área 
     sobre los bloques lógicos."
\section{Objetivos}

\subsection{General}
            Evaluar el impacto de la variación de parámetros estructurales en el 
            rendimiento, área y potencia de una arquitectura abierta de lógica 
            reconfigurable a nivel de simulación y pruebas físicas utilizando una 
            FPGA comercial como plataforma de emulación física.
\subsection{Específicos}
	\begin{itemize}
                \item Seleccionar una arquitectura de lógica reconfigurable de hardware 
                abierto y su flujo de herramientas asociado, justificando la elección 
                basado en su nivel de parametrización, soporte de verificación funcional 
                e historial de implementación física en otros proyectos.

                \item Configurar e instanciar la arquitectura seleccionada en una tarjeta 
                FPGA y caracterizarla basandose en reportes de síntesis, ruteo y ejecución 
                funcional, además de implementar al menos dos circuitos de prueba en dicha 
                lógica programable.
                             
                \item Caracterizar el impacto de la variación de los parámetros 
                estructurales disponibles en la arquitectura de lógica programable sobre 
                la utilización de recursos, el consumo de potencia y el rendimiento a nivel 
                de temporización.
    \end{itemize}
 
\section{Lista de Tareas por Objetivo}

Esta sección se enfoca en proporcionar un desglose detallado de las tareas que 
deben realizarse para lograr cada objetivo específico en el proyecto. Esta 
estructura por objetivo facilita la comprensión de las tareas individuales en el 
contexto de los objetivos más amplios del proyecto. Cada objetivo del proyecto 
tendrá su propia lista de tareas relacionadas, lo que proporciona una visión clara 
de lo que se debe lograr para alcanzar ese objetivo. Estas tareas se presentan de 
manera secuencial, lo que indica el orden en el que deben completarse.

\textbf{\normalsize{Objetivo Específico 1}}

Este objetivo se centra en la investigación de la arquitectura de 
lógica reconfigurable y el flujo de harramientas para su implementación, 
considerando la el soporte para su parametrización, verificación 
funcional e historial de implementación física.

\begin{enumerate}

    \item Investigar y recopilar información sobre las distintas 
    arquitecturas de lógica reconfigurable de hardware abierto 
    disponibles.
     
    \item Evaluar la flexibilidad y el nivel de parametrización 
    estructural que ofrece cada arquitectura preseleccionada.
    
    \item Analizar la madurez del flujo de herramientas asociado, 
    comprobando la compatibilidad y el soporte para la verificación 
    funcional.

    \item Revisar el historial de implementación física en silicio 
    de las opciones evaluadas para constatar su viabilidad y 
    confiabilidad en proyectos previos. 

    \item Seleccionar la arquitectura definitiva y justificar técnicamente 
    la elección, basada en el cumplimiento de los criterios previamente 
    analizados.

\end{enumerate}

Las tareas descritas justifican el uso de la arquitectura en las próximas 
etapas y su implementación en la FPGA.

\textbf{\normalsize{Objetivo Específico 2}}

El segundo objetivo se centra en la configuración de la arquitectura 
seleccionada mediante JTAG y la caracterización basada en reportes 
desplegados por las herramientas seleccionadas en la etapa anterior, 
además de la implementación de dos circuitos de prueba.

\begin{enumerate}

    \item Ajustar los módulos de la arquitectura seleccionada para asegurar 
    su compatibilidad con la herramienta de síntesis de la FPGA objetivo.
    
    \item Integrar la arquitectura en el proyecto principal de la herramienta de desarrollo.
    
    \item Correr el proceso de síntesis y extraer los reportes de la cantidad de LUTs, 
    flip-flops y bloques de memoria requeridos.

    \item Ejecutar la implementación en la FPGA y compilar los reportes de tiempos 
     de setup y hold, consumo de potencia estimado y congestión de ruteo.

    \item Generar el bitstream necesario para configurar la FPGA física.
    
    \item Seleccionar dos circuitos que estresen diferentes aspectos de la lógica programable.
    
    \item Codificar ambos circuitos independientemente.
    
    \item Escribir los testbenches correspondientes y simular el comportamiento 
    de los circuitos de prueba para asegurar su corrección funcional antes de la 
    implementación física.

    \item Compilar y cargar el primer y segundo circuito de prueba sobre la 
    arquitectura de lógica programable ya instanciada.

    \item Estimular las entradas en el hardware real, capturar las salidas y contrastar la 
    ejecución funcional contra los resultados obtenidos en las simulaciones

\end{enumerate}

Estas tareas permiten validar la arquitectura seleccionada y dan información benéfica para 
el desarrollo de la caracterización de la arquitectura seleccionada.

\textbf{\normalsize{Objetivo Específico 3}}

El tercer objetivo da base a la caracterizacion de la arquitecura de lógica programable 
seleccionada según la utilización de recursos.

\begin{enumerate}

    \item Identificar cuáles variables de la arquitectura de la arquitectura lógica 
    programable serán modificables.

    \item Definir un conjunto discreto de valores para cada parámetro seleccionado y 
    limitar las combinaciones posibles a un número manejable y representativo.
    
    \item Recopilar la cantidad de bloques lógicos básicos utilizados y la estimación de 
    área equivalente de silicio.

    \item Extraer del Análisis de Temporización Estática y el retardo de la ruta crítica.

    \item Compilar las estimaciones de consumo de potencia dinámica y estática reportadas 
    tras el ruteo.

    \item Procesar los datos consolidados para graficar el comportamiento de las métricas 
    frente a la variación de cada parámetro.

    \item Analizar los puntos donde la mejora de una métrica perjudica a otra (Trade-off).
    
    \item Concluir cuáles son las configuraciones de parámetros estructurales que ofrecen 
    el mejor balance técnico para los casos de uso proyectados.

\end{enumerate}


\section{Alcances y Limitaciones del Proyecto}

En la sección de limitaciones se exponen los desafíos y restricciones propios de la 
naturaleza del proyecto. Si bien estas condiciones pueden representar obstáculos, 
al mismo tiempo delimitan el alcance del trabajo y establecen un marco de referencia 
realista para su ejecución. A continuación, se detallan las principales limitaciones 
identificadas:

\begin{itemize}

    \item La ejecución de los experimentos se llevará a cabo con la FPGA 
    disponibles durante la evaluación del proyecto. Esta limitación en los recursos 
    físicos para la realización de pruebas podría restringir tanto la diversidad de 
    los escenarios evaluados como la naturaleza de los datos susceptibles de recolección.

    \item Al instanciar una arquitectura de lógica programable sobre una FPGA comercial, 
    el área, la potencia y el retardo de temporización estarán limitados por la arquitectura 
    de la FPGA anfitriona.

    \item La configuración óptima identificada estará restringida a las características 
    de los circuitos seleccionados y podría no ser generalizable a todas las aplicaciones 
    digitales.

    \item No se realizará una exploración exhaustiva de todo el espacio de diseño, por lo 
    que la configuración óptima que maximiza el rendimiento corresponderá únicamente a un 
    óptimo dentro del subconjunto explorado.

    \item Las métricas de consumo de potencia se extraerán a partir de reportes generados 
    por herramientas de síntesis y ruteo, no mediante mediciones físicas directas sobre 
    la arquitectura implementada por lo tanto se introduce un margen de error inherente a 
    los modelos matemáticos y estimaciones.



\end{itemize}